\id{IRSTI 65.59.03}{}

\begin{header}
\swa{A.E.~Abdurasulova, B.A.~Rskeldiev}{SHEEP MEAT AND PLANT ANTIOXIDANTS IN THE TECHNOLOGY OF HIGH NUTRITIONAL-VALUE MEAT ROULADE: A SYSTEMATIC REVIEW}

A.E.~Abdurasulova\envelope,
B.A.~Rskeldiev
\end{header}

\begin{affil}
Almaty Technological University, Almaty, Kazakhstan

\corrauthor{Corresponding author: abdurasulova.aigerim@mail.ru}
\end{affil}

The relevance of this work is due to the need to develop functional meat
products with increased nutritional value that are also resistant to
oxidative degradation, which is exacerbated in red meat due to the
presence of heme pigments and reactive lipid fractions. The research
hypothesis is that combining sheep meat with the bioactive components of
sea buckthorn and grape seeds will create a biochemically sound
antioxidant platform for meat roulade, reducing oxidative losses and
stabilising quality during storage. The aim of the article is to
critically synthesise current data on the nutritional value and
technological characteristics of lamb, as well as on the composition and
functional effects of sea buckthorn and grape seeds in meat systems. The
methods included a systematic review of publications from 2015 to 2025,
searching in English, Russian, and Kazakh and selecting works containing
analytical quality indicators. The results show that the bioactive
fractions of sea buckthorn (vitamin C at around 300-600 mg/100 g,
carotenoids at 5--36 mg/100 g, as well as phenolic compounds) and the
polyphenolic complex of grape seeds (proanthocyanidins - 50-200 mg/g in
extracts, flavanols and dietary fibre) are most consistently associated
in studies with a reduction in oxidation markers, more stable colour
characteristics and a variable, dose-dependent effect on textural and
technological indicators. The most controllable form of application for
roll systems appears to be extracts, which ensure accurate dosing and
uniform distribution of bioactive components. Future studies will be
based on optimising the amount of plant antioxidants added and the
stages of addition, heat treatment and packaging, while controlling the
trade-offs between antioxidant effect and organoleptic characteristics.

{\bfseries Keywords:} sheep meat, meat roulade, functional meat products,
sea buckthorn, grape seeds, oxidative stability.

\begin{header}
БАРАНИНА И РАСТИТЕЛЬНЫЕ АНТИОКСИДАНТЫ В ТЕХНОЛОГИИ МЯСНОГО РУЛЕТА ПОВЫШЕННОЙ ПИЩЕВОЙ ЦЕННОСТИ: СИСТЕМНЫЙ ОБЗОР

А.Е.~Абдурасулова\envelope,
Б.А.~Рскелдиев
\end{header}

\begin{affil}
Алматинский технологический университет, г. Алматы, Казахстан,

e-mail: abdurasulova.aigerim@mail.ru
\end{affil}

Актуальность работы обусловлена необходимостью разработки функциональных
мясных продуктов с повышенной пищевой ценностью и одновременно
устойчивой к окислительной деградации, которая в красном мясе
усиливается из-за присутствия гемовых пигментов и реакционноспособных
липидных фракций. Гипотеза исследования состоит в том, что сочетание
баранины с биоактивными компонентами облепихи и виноградных косточек
позволит сформировать биохимически обоснованную антиоксидантную
платформу для мясного рулета, обеспечивающую снижение окислительных
потерь и стабилизацию качества при хранении. Цель статьи --- критически
синтезировать современные данные о пищевой ценности и технологических
особенностях баранины, а также о составе и функциональных эффектах
облепихи и виноградных косточек в мясных системах. Методы включали
системный обзорный анализ публикаций 2015--2025 гг. с поиском на
английском, русском и казахском языках и отбором работ, содержащих
аналитические показатели качества. Результаты показывают, что
биоактивные фракции облепихи (витамин C порядка 300-600 мг/100 г,
каротиноиды 5--36 мг/100 г, а также фенольные соединения) и
полифенольный комплекс виноградных косточек (проантоцианидины, флаванолы
и пищевые волокна) в исследованиях наиболее устойчиво ассоциируются со
снижением маркеров окисления, более стабильными характеристиками цвета и
вариабельным, дозозависимым влиянием на текстурно-технологические
показатели. Наиболее управляемой формой внесения для рулетных систем
представляются экстракты, обеспечивающие точную дозировку и равномерное
распределение биоактивных компонентов. Будущие исследования будут
основаны на оптимизации количество вносимых растительных антиоксидантов
и на этапах внесения, режима термообработки и упаковки при контроле
компромиссов антиоксидантный эффект -- органолептические характеристики.

{\bfseries Ключевые слова:} баранина, мясной рулет, функциональные мясные
продукты, облепиха, виноградные косточки, окислительная стабильность.

\begin{header}
ҚОЙ ЕТІ ЖӘНЕ ӨСІМДІК ТЕКТІ АНТИОКСИДАНТТАР ЖОҒАРЫ ТАҒАМДЫҚ ҚҰНДЫЛЫҒЫ БАР ЕТ РУЛЕТІНІҢ ТЕХНОЛОГИЯСЫНДА: ЖҮЙЕЛІ ШОЛУ

А.Е.~Абдурасулова\envelope,
Б.А.~Рскелдиев
\end{header}

\begin{affil}
Алматы технологиялық университеті, Алматы, Қазақстан,

e-mail: abdurasulova.aigerim@mail.ru
\end{affil}

Жұмыстың өзектілігі тағамдық құндылығы жоғары және сонымен қатар тотығу
деградациясына төзімді функционалды ет өнімдерін әзірлеу қажеттілігімен
байланысты, өйткені қызыл ет құрамындағы гем пигменттері мен реакцияға
қабілетті липидтік фракциялар тотығу үдерістерін күшейтеді. Зерттеу
гипотезасы қой етін шырғанақ пен жүзім сүйектерінің биоактивті
компоненттерімен біріктіру ет рулетіне арналған биохимиялық тұрғыдан
негізделген антиоксиданттық платформаны қалыптастырып, тотығу шығындарын
азайтуға және сақтау кезінде сапаны тұрақтандыруға мүмкіндік береді
деген тұжырымға негізделеді. Мақаланың мақсаты --- қой етінің тағамдық
құндылығы мен технологиялық сипаттамалары, сондай-ақ облепиха мен жүзім
сүйектерінің құрамы және олардың ет жүйелеріндегі функционалдық әсерлері
жөніндегі заманауи деректерді сыни тұрғыдан жинақтау. Зерттеу әдістері
2015-2025 жылдар аралығындағы жарияланымдарға жүйелі шолу жүргізуді,
ағылшын, орыс және қазақ тілдерінде деректер іздеуді және сапаның
аналитикалық көрсеткіштері келтірілген еңбектерді іріктеуді қамтыды.
Алынған нәтижелер шырғанақтың биоактивті фракциялары (С дәрумені шамамен
300-600 mg/100 g, каротиноидтар 5-36 mg/100 g, сондай-ақ фенолдық
қосылыстар) және жүзім сүйектерінің полифенолдық кешені
(проантоцианидиндер, флаванолдар және тағамдық талшықтар) зерттеулерде
тотығу маркерлерінің төмендеуімен, түс көрсеткіштерінің неғұрлым тұрақты
болуымен және құрылымдық-технологиялық сипаттамаларға дозасына тәуелді
өзгермелі әсерімен ең тұрақты түрде байланысты екенін көрсетті. Рулет
типті жүйелер үшін биоактивті компоненттердің дәл мөлшерленуін және
біркелкі таралуын қамтамасыз ететін енгізу формасы ретінде экстрактылар
ең басқарылатын нұсқа болып табылады. Болашақ зерттеулер өсімдік текті
антиоксиданттардың енгізілетін мөлшерін, енгізу кезеңдерін, термиялық
өңдеу мен қаптау режимдерін оңтайландыруға, сондай-ақ антиоксиданттық
әсер мен органолептикалық қасиеттер арасындағы өзара ымыраны бақылауға
бағытталатын болады.

{\bfseries Түйін сөздер:} қой еті, ет рулеті, функционалды ет өнімдері,
шырғанақ, жүзім сүйектері, тотығу тұрақтылығы..

\begin{multicols}{2}
{\bfseries Introduction.} In modern meat science, the concept of functional
meat products has shifted from simple fortification toward
nutrient-dense, technologically stable matrices that can retain quality
during processing and storage. A central scientific problem is that many
of the nutrients that make meat valuable coexist with pro-oxidant
drivers (heme iron, unsaturated lipids, oxygen exposure during
comminution), which accelerate lipid oxidation, impair colour stability,
and generate off-flavours which ultimately shortening shelf life and
reducing consumer acceptance {[}1{]}.

Sheep meat is often positioned as a nutrient-dense red meat because it
supplies complete proteins and a spectrum of micronutrients. In a recent
lamb loin dataset, Kenenbay et al. reported iron around 2,0-2,2 mg, zinc
around 3,0-3,3 mg, vitamin E around 1,0-1,4 mg, and vitamin
B\tsb{12} around 0,50--0,60 mg, illustrating the micronutrient
``core'' that motivates lamb-based functional products {[}2{]}. At the
same time, sheep meat can be technologically vulnerable because
intramuscular fat composition includes a meaningful PUFA fraction, which
supplies oxidation substrates. In the same work, PUFA accounted for
roughly 21,8-23,5\% of total fatty acids in intramuscular fat, a range
that is nutritionally attractive but chemically reactive in oxygenated,
iron-rich matrices {[}3{]}.

Sea buckthorn (\emph{Hippophae rhamnoides}) is scientifically
interesting for meat applications because it combines multiple
antioxidant compounds. Sea buckthorn contains phenolics-radical
scavenging and metal chelation, carotenoids-quenching of reactive
species and contribution to colour tone, and vitamin C-redox cycling and
potential support of metmyoglobin reduction depending on dose and matrix
conditions. In a recent compositional synthesis, Olas et al. highlighted
sea buckthorn berries as a notably vitamin-rich raw material, reporting
vitamin C values that can reach approximately 275 mg/100 g, alongside
lipid-soluble antioxidants such as vitamin E around 3,54 mg/100 g
{[}4{]}. Mechanistically, these compounds may influence meat quality
through several routes - termination of lipid peroxyl radicals by
phenolic hydrogen donation, chelating Fe³⁺, lowering catalytic redox
cycling that accelerates peroxide decomposition and modulating pH via
organic acids {[}5{]}.

Grape seeds represent a high-value secondary resource from viticulture
and wine processing, aligning with circular-economy logic while offering
dense polyphenolic chemistry. A recent valorisation review reported that
grape by-products commonly exhibit high total phenolic content, with
representative ranges for different fractions often reaching the
hundreds of mg GAE/100 g, supporting their use as antioxidant carriers
in foods. The dominant antioxidant ``identity'' of grape seed materials
is typically associated with proanthocyanidins (condensed tannins) -
polymeric flavanols that can suppress oxidation by radical scavenging,
metal chelation, and interruption of chain propagation. Verma et al.
summarised proanthocyanidins as oligopolymeric flavanol structures with
broad bioactivity and technological relevance, which is directly
applicable to red meat systems where oxidation is often the limiting
factor for quality stability {[}6, 7{]}.

Given mutton's nutrient density and red-meat bioavailability advantages,
its susceptibility to oxidative and colour degradation, and the
complementary antioxidant profiles of sea buckthorn and grape seeds, we
hypothesise that a mutton-based meat roulade can be strengthened
nutritionally while improving oxidative and quality stability through
targeted inclusion of these plant resources.

According to the hypothesis, the aim of this systematic review is to map
and critically synthesise contemporary evidence on lamb meat quality
constraints relevant to processing and storage, sea buckthorn-derived
ingredients as multifunctional antioxidant inputs, and grape seed
ingredients as polyphenol-forward stabilisers and framing a biochemical
rationale for their combined use in nutrient-dense lamb product
development.

{\bfseries Materials and methods.} A structured literature search for
systematic review was performed in Scopus, Web of Science Core
Collection, PubMed, and ScienceDirect, and the reference lists of
influential papers addressing sheep meat quality, safety, and storage
were additionally screened to identify other relevant studies. To
strengthen regional representation, peer-reviewed articles published in
Kazakhstani journals listed by the Science and Higher Education Quality
Assurance Committee of the Ministry of Science and Higher Education of
the Republic of Kazakhstan as recommended outlets for reporting primary
scientific results were also included. The search was limited to the
last 10 years.

To minimize language bias and capture locally specific evidence,
searches were conducted in English, Russian, and Kazakh. The search
strategy combined terms describing the raw material, the product
concept, the plant additives, and outcome domains. Variants of «sheep
meat» or «lamb» or «mutton» were used alone and also together with terms
describing the target product («meat roulade», «rolled meat»,
«restructured meat product», «meat roll») and the additives («sea
buckthorn» or «Hippophae rhamnoides» and «grape seed» or «grape seed
extract»), along with keywords representing the outcome domains with
equivalent terminology in Russian and Kazakh. Boolean logic, truncation,
and platform-specific filters were adjusted to each database.

Studies were excluded if they used matrices irrelevant to meat products
or did not employ a meat model system, as well as works that lacked
analytical assessment of meat quality, including physicochemical,
oxidative, colour, texture, microbiological, or sensory parameters.
Studies that were not peer‑reviewed full‑text articles, such as
abstracts without complete data, or that did not provide sufficient
methodological detail for interpretation, were also excluded. In
addition, research focused on additives unrelated to the scope of this
review, that is, without a clear connection to sea buckthorn bioactives
or grape‑seed‑derived phenolic fractions, was omitted unless such
studies were cited only to provide contextual background.

Records retrieved from databases were compiled and deduplicated. Titles
and abstracts were screened for relevance to ovine/red-meat systems, sea
buckthorn or grape seed--derived ingredients, and quality outcomes. Full
texts were then assessed against the eligibility criteria. From each
included study, data were extracted on meat matrix and product type,
additive form and preparation (extract, powder, pomace, oil-associated
fraction), dose level and incorporation stage, and outcome measures.

Preference in interpretation was given to studies with clearly described
processing conditions, validated analytical methods, and storage trials
reflecting realistic shelf-life scenarios. Where findings were
inconsistent, differences were interpreted in relation to matrix
composition, dose, incorporation stage, and thermal/storage conditions
rather than treating outcomes as directly comparable.

{\bfseries Results and discussion.} Multiple independent studies
consistently confirm that lamb is a nutrient-dense red meat, providing
high-quality protein together with bioavailable micronutrients.
Junkuszew et al. analysed sheep breeds and reported a dry matter range
of 24,11-27,13\%, which implies moisture 72,87-75,89 g/100 g. In the
same muscles, intramuscular fat was 4,60-5,50 g/100 g, total protein was
19,95-22,18 g/100 g, and ash were 0,71-1,06 g/100 g, showing that lean
lamb muscle is primarily a high-water, high-protein matrix with moderate
intramuscular lipid {[}8{]}. At the biochemical level, lamb's protein
fraction is characterised by a dense pool of essential amino acids
(EAAs) that support both human nutrition and meat functionality -- water
binding, gelation, texture formation. Total essential amino acids of
24,17-27,20 g/100 g (dry matter) and non-essential amino acids of
35,13-39,48 g/100 g (dry matter). Among abundant contributors were
glutamic acid 9,03-9,62, leucine 4,31-4,80, and lysine 3,69-4,17 g/100 g
(dry matter). Such an amino acid distribution is relevant not only
nutritionally but also sensorially, because glutamate-associated
pathways contribute to savoury perception after proteolysis {[}9{]}.

From a lipid biochemistry perspective, lamb's nutritional profile is
strongly shaped by ruminant lipid metabolism (biohydrogenation and
isomerisation in the rumen) and by feeding system. In lamb Longissimus
dorsi, Sabbioni et al. reported PUFA 25\% of total fatty acids, a
PUFA/SFA ratio of 0,68, with n-3 = 8\% and n-6 =14\%, alongside
favourable composite lipid indices- atherogenic index = 0,56-0,67 and
thrombogenic index = 0,78-0,96. Biochemically, higher PUFA and n-3
fractions increase membrane fluidity and nutritional value, but also
raise susceptibility to lipid peroxidation-so oxidative stability
becomes a key quality dimension in lamb during storage and cooking
{[}10{]}.

All studies included in the tables underwent screening and were selected
based on the inclusion/exclusion criteria specified in Materials and
Methods.
\end{multicols}

\tcap{Table 1 - Comparative chemical composition of lamb, beef and pork meat}
\begin{longtblr}[
  label = none,
  entry = none,
]{
  cells = {c},
  cells = {font = \small},
  row{1} = {font=\bfseries\small},
  hlines,
  vlines,
}
Parameters & Mutton & Source & Beef & Source & Pork & Source\\
Moisture, g/100 g & 68,1 & {[}11] & 72,4 & {[}13] & 74,9 & {[}16]\\
Protein, g/100 g & 18,5 & {[}11] & 19,1 & {[}13] & 18,2 & {[}16]\\
Fat, g/100 g & 12,1 & {[}11] & 12,27 & {[}13] & 12,34 & {[}16]\\
Ash, g/100 g & 0,99 & {[}11] & 1,2 & {[}13] & 1,07 & {[}16]\\
Calcium, mg/100 g & 7,27 & {[}11] & 4,39 & {[}14] & 10 & {[}16]\\
Iron, mg/100 g & 1,6–2,4 & {[}11] & 1,39 & {[}14] & 0,8 & {[}16]\\
Potassium, mg/100 g & 301 & {[}11] & 367 & {[}14] & 330 & {[}16]\\
Magnesium, mg/100 g & 20,8 & {[}11] & 21,3 & {[}14] & 20,4 & {[}16]\\
Sodium, mg/100 g & 67,5 & {[}11] & 40,4 & {[}14] & 50,1 & {[}16]\\
Phosphorus, mg/100 g & 193 & {[}11] & 163 & {[}14] & 220 & {[}16]\\
Zinc, mg/100 g & 3,08 & {[}11] & 5,28 & {[}14] & 2,4 & {[}16]\\
Selenium, нg/100 g & 12,3 & {[}11] & 12,5 & {[}14] & 21 & {[}16]\\
Thiamine (B1), mg/100 g & 0,14 & {[}12] & 0,06 & {[}15] & 0,72 & {[}17]\\
Riboflavin (B2), mg/100 g & 0,07 & {[}12] & 0,12 & {[}15] & 0,21 & {[}17]\\
Niacin (B3), mg/100 g & 3,32 & {[}12] & 2,4 & {[}15] & 4,7 & {[}17]\\
Pantothenic acid (B5), mg/100 g & 0,15 & {[}12] & 0,12 & {[}15] & 0,47 & {[}17]\\
Pyridoxine (B6), mg/100 g & 0,54 & {[}12] & 0,2 & {[}15] & 0,72 & {[}17]\\
Folate (B9), mkg/100 g & 4,16 & {[}12] & 5,9 & {[}15] & 5,2 & {[}17]\\
Cyanocobalamin (B12), mkg/100 g & 2,53 & {[}12] & 0,61 & {[}15] & 1,12 & {[}17]\\
α-Tocopherol (E), mg/100 g & 0,31 & {[}12] & 0,21 & {[}15] & 0,26 & {[}17]
\end{longtblr}

\begin{multicols}{2}
A consistent biochemical advantage of lamb over pork is iron density,
which matters because iron in meat is largely present as heme iron
(higher bioavailability than non-heme forms). Junkuszew et al. reported
lamb iron 2,60-2,93 mg/100 g, while Aaslyng et al. report lean pork loin
iron 1,6--2,4 mg/100 g and beef loin iron 2,40 mg/100 g. Practically,
this places lamb far above pork and at least comparable to beef in total
iron, supporting lamb as a strong option when dietary iron adequacy is a
priority {[}18{]}.

A more ``technological'' advantage of lamb versus pork is that its
matrix naturally combines moderate intramuscular fat with a high EAA
density, which supports both flavour development and a firm protein
network after cooking. Intramuscular fat 4,60-5,50 g/100 g and with
EAA/NEAA ratio = 0,68-0,72, a balance that helps explain why lamb can
maintain sensory richness without requiring high fat levels typical of
some pork products. Compared with very lean beef (2,27 g fat/100 g),
lamb's higher intramuscular lipid fraction can improve perceived
juiciness and flavour intensity-while still staying within a moderate
fat range {[}19{]}.

Sea buckthorn (\emph{Hippophae rhamnoides}) is biochemically unusual
among edible berries because its core matrix combines high acidity,
moderate sugars and very high levels of redox-active micronutrients. In
a compositional survey of sea buckthorn fruits, Tkacz et al. reported
total soluble solids around 3,38-5,67 g/100 g, with glucose 1,82-3,13
g/100 g and fructose 0,90-1,71 g/100 g, while the sum of organic acids
reached 0,97-2,80 g/100 g. This sugar-acid architecture explains both
the characteristic sharp taste and the technological value of sea
buckthorn as a naturally acidic, antioxidant-rich ingredient {[}20{]}.

A key nutritional hallmark is ascorbic acid (vitamin C), which in sea
buckthorn can reach physiologically meaningful levels for a realistic
serving size. Bośko et al. quantified ascorbic acid at 300-600
mg/100 g of fruit, a 50-100 g portion can plausibly deliver tens to less
than 100 mg vitamin C {[}21{]}. Biochemically, ascorbate acts as a
two-electron reductant supporting the cellular antioxidant network, and
it also serves as an enzymatic cofactor (e.g., hydroxylation reactions
in connective-tissue metabolism). Importantly for a review, reported
vitamin C values vary strongly with genotype and harvest conditions;
Jaroszewska et al. summarize ranges where measured vitamin C spans
roughly 61-159 mg/100 g, reinforcing that region must be treated as a
major effect modifier {[}22{]}.

Beyond vitamins, sea buckthorn is rich in polyphenols, especially
flavonol derivatives, which matter because they participate in radical
scavenging, metal chelation, and redox-signalling modulation. In berry
extracts from multiple cultivars, Criste et al. reported total phenolics
10,12-18,66 mg GAE/g and flavonoids 6,57-9,01 mg QE/g, showing
cultivar-dependent but consistently high phenolic density for an edible
fruit {[}23{]}. From a human-nutrition angle, these phenolics are
relevant not only as direct antioxidants, but also because many act as
substrates for gut microbial metabolism, generating smaller phenolic
acids that can influence inflammatory and metabolic signalling.

Sea buckthorn's carotenoid system adds a second, lipophilic antioxidant
layer with provitamin-A potential. Zenkova et al. showed that total
carotenoids differed sharply by variety (5,63 to 35,78 mg/100 g), and
they identified zeaxanthin as a predominant compound, reaching 27,78
mg/100 g in one cultivar, lutein up to 4,74 mg/100 g and β-carotene
around 0,17-1,94 mg/100 g {[}24{]}. Functionally, carotenoids are
efficient physical quenchers of singlet oxygen and interceptors of lipid
radicals in membranes and lipoproteins-mechanisms that complement
vitamin C's aqueous redox role and help explain why sea buckthorn is
repeatedly positioned as a multi-compartment antioxidant food.

In a controlled dietary intervention, Huang et al. showed that
increasing sea buckthorn oil intake raised circulating trans-palmitoleic
acid by 26,6\%, confirming that sea buckthorn oils measurably alter
circulating lipid markers in humans {[}25{]}.

Sea buckthorn contributes minerals and may support metabolic outcomes
through combined fiber-polyphenol-lipid mechanisms. In compositional
measurements, Kim et al. reported high potassium 600,918 mg/100 g,
with notable calcium 169,546 mg/100 g, magnesium 120,321 mg/100 g,
phosphorus 135,445 mg/100 g, and iron 3,203 mg/100 g {[}26{]}. On the
functional side, a randomized crossover trial by Ren et al. in
individuals with impaired glucose regulation found that after sea
buckthorn fruit puree, fasting plasma glucose changed from 6,04 to
5,80 mmol/L, and the between-condition comparison for fasting plasma
glucose was significant (p = 0,045), suggesting a modest but
detectable glycemic signal in a real dietary context {[}27{]}.

In meat systems, sea buckthorn is typically introduced through two
practical formats, each targeting a slightly different set of mechanisms
(Table 2). First, water-compatible extracts or juices are incorporated
directly into the comminuted phase during mixing, so the bioactives
become distributed throughout the protein-fat matrix. Second, dry
powders (extract powders or pomace powders) are added at the dry
blending stage or during final mixing before stuffing or forming. This
route increases the local concentration of solids and can change
hydration dynamics.

Across studies, the most consistently affected outcomes are markers of
oxidative stability, because sea buckthorn supplies multiple
redox-active pools. In biochemical terms, lipid oxidation in meat is a
chain process initiated by reactive oxygen species and catalysed by heme
iron and other transition metals. It propagates through formation of
lipid radicals (L•), peroxyl radicals (LOO•) and hydroperoxides (LOOH),
eventually yielding secondary aldehydes measured by TBARS. Sea buckthorn
phenolics can donate electrons or hydrogen atoms to terminate LOO•
radicals (chain-breaking antioxidant action), while also chelating
Fe²⁺/Fe³⁺, reducing Fenton-type initiation. At the same time,
carotenoids can quench singlet oxygen and intercept radicals in lipid
regions, and ascorbate can regenerate oxidized antioxidant
forms-although its role is dose-and matrix-dependent because in the
presence of free iron, strong reducers can sometimes show pro-oxidant
behaviour. Practically, this is why sea buckthorn additions frequently
translate into lower TBARS, POV, AV and improved colour stability in
chilled storage, especially in systems with higher fat or higher heme
pigment load.
\end{multicols}

\tcap{Table 2 - Application of sea buckthorn and its effects in meat systems}
\begin{longtblr}[
  label = none,
  entry = none,
]{
  width = \linewidth,
  colspec = {Q[110]Q[150]Q[187]Q[446]Q[60]},
  cells = {font = \small},
  row{1} = {c,font=\bfseries\small},
  column{5} = {c},
  hlines,
  vlines,
}
Meat or product type & Sea buckthorn form & Dose and incorporation stage & Results & Source\\
Horse meat, chunked delicacy “Jaya” & Dry sea buckthorn extract (powder) & 3\% (w/w), mixed into formulation prior to processing & TBARS reduced by 26,5–44,7\%, remaining 0,26–0,49 mg MDA/kg; AV 0,28–0,59 mg KOH/g; POV 0,11–0,55 meq O₂/kg; hardness decreased 1,41–1,33; & {[}28]\\
Cooked-smoked sausages & Sea buckthorn extract (with hemp protein as co-ingredient) & 0,005\% – 0,015\% added during batter formulation before stuffing & Best sensory acceptance reported at 0,01\%. Total microbial counts within acceptable limits up to 36 days, supporting shelf life under 0–4 °C & {[}29]\\
Pork sausages & Sea buckthorn extract & 3 mL/kg and 5 mL/kg, incorporated during mixing & pH: 6,55(control) and 6,18 (5 mL/kg); colour day: L* 67,91 (control), 66,54 (3 mL/kg) and 66,28 (5 mL/kg); Firmness: 385,96 N (control), 371,77 N (3 mL/kg) and 407,99 N (5 mL/kg); & {[}30]\\
Beef burgers & Sea buckthorn berry juice & Juice incorporated at the patty formulation stage (before cooking) & Cooking loss decreased from 41,2\% (control) to 38,9\%. Yield increased from 58,8\% to 62,0\%. Antioxidant activity in burgers reported in the range 42–440 µmol/L/100 g. & {[}31]\\
Chilled beef under super-chilled storage & Sea buckthorn pomace extract & SBP extract 0\% / 1\% / 2\% / 3\% (w/w) & Shelf life prolonged from 25 to 35 days with 3\% extract film. Higher L* and a* maintained. & {[}32]\\
Pork sausages & Sea buckthorn extract & Extract used at 3 mL/kg and 5 mL/kg (incorporation into sausage system) & Extract characterised by TPC 677,58 g GAE/kg and DPPH inhibition 65,04\%. Authors report reduced malondialdehyde (MDA) formation in sausages (oxidation), while antimicrobial impact was not evident & {[}33]
\end{longtblr}

\begin{multicols}{2}
Grape seeds (\emph{Vitis vinifera L.)} are nutritionally meaningful
mainly because they concentrate dietary fibre and plant secondary
metabolites. In a compositional study of defatted grape seed meal,
Baca-Bocanegra \emph{et al.} reported moisture 11,78\%, ash 2,87\%,
dietary fibre 22,86\%, protein 10,76\%, lipid 1,30\%, and available
carbohydrates 50,43\%. This profile is typical of an ingredient where
nutritional value is not driven by fat or simple sugars, but by
fibre-associated physiological effects and by co-localised phenolics
bound to the cell-wall matrix {[}34{]}.
\end{multicols}

\tcap{Table 3 - Application of grape seed and its effects in meat systems}
\begin{longtblr}[
  label = none,
  entry = none,
]{
  colspec = {Q[110]Q[100]Q[250]Q[350]Q[60]},
  cells = {font = \small},
  row{1} = {c,font=\bfseries\small},
  column{5} = {c},
  hlines,
  vlines,
}
Meat or product type & Form of grape seed ingredient & Dose and incorporation method & Results & Source\\
Pig liver pâté, heat-treated, & Grape seed extract & 50 / 200 / 1000 mg/kg, incorporated into the meat batter during mixing & Baseline TBARS 0,4–3,3 mg MDA/kg; higher extract level resulted in lower oxidative deterioration. & {[}38]\\
Dry-cured pork sausage (“chorizo”) & Grape seed extract & Extract dose range - 0,05–1,0 g/kg; incorporated during vacuum mincing & TBARS  0,4 mg MDA/kg; hardness decreased, with minimum values reported at 0,20 g/kg. TVC in extract-treated batches up to 7,88–7,97 log CFU/g. & {[}39]\\
Ground beef, retail display storage & Grape seed extract & 250 mg/kg, added during minced meat mixing & TBARS (control) - 0,96 mg MDA/kg on day 6; with extract - 0,82 mg MDA/kg on day 6. & {[}40]\\
Meat patties (85\% beef + 15\% pork fat) & Grape seed extract & 0,10 / 0,25 / 0,50 / 0,75 g/kg; extract pre-dissolved in water and incorporated during mixing before patty forming & Colour on day 10 a* = 8,62 (control) and 13,20 (0,75 g/kg). TBARS maintained within 0–0,51 mg MDA/kg. & {[}41]\\
Dry-cured pork neck, long-term storage & Grape seed extract & Dose variants GSE-1 / GSE-2 / GSE-5. Incorporated into the meat mass at the beginning of processing & After 2 months TBARS 0,72 mg MDA/kg (control) and 0,42 mg MDA/kg (GSE-5), ascorbic acid - 0,53 mg MDA/kg. & {[}42]\\
Beef–pork sausage model system & Grape seed extract & 0,125 / 0,25 g per 250 g meat batter added during mixing. Some treatments combined with ascorbic acid & TBARS shows 0,43 mg MDA/kg (GSE 0,125 g/250 g), 0,41 mg MDA/kg (GSE + AA); and 0,45 mg MDA/kg (control). & {[}43]\\
Smoked pork sausage & Grape seed extract & Extract incorporated during mixing & Reported reductions TBARS by 30,3\%, expressible water by 26,4\%. Sensory score higher by 18,0\%; TVC moulds also decreased. & {[}44]
\end{longtblr}

\begin{multicols}{2}
The lipid fraction of grape seeds is also biochemically relevant when
the oil is recovered. Carmona-Jiménez \emph{et al.} quantified grape
seed oil yields of 12,46-22,52 g oil/100 g dry seed and showed that the
oil is strongly skewed toward unsaturated fatty acids (USFA 82-84\%),
with PUFA 64,04-69,44\% and MUFA 16,43-20,27\%. Linoleic acid dominated
at 65,55-69,00\%, while oleic acid ranged up to about 20,13\%. In the
same research line, grape seed oils were reported to contain total
tocopherols 196,2-248,0 mg/kg oil, supporting Vitamin-E-linked
antioxidant functionality of the lipid phase {[}35{]}. Ivanov \emph{et
al.} measured high concentrations in grape seed-derived ingredients,
reporting total phenolic content 111,22 mg GAE/g, total flavonoids 51,50
mg QE/g, and total proanthocyanidins 170,45 mg CE/g {[}36{]}. These
compounds are rich in pyrogallol hydroxyl groups, which enables radical
scavenging via electron transfer, and metal chelation, which can slow
Fenton-type propagation of oxidation in biological and food systems.

Human evidence (mostly from grape seed extracts) suggests measurable
effects on vascular and redox endpoints. In a meta-analysis of
randomised controlled trials, Foshati \emph{et al.} found that grape
seed extract supplementation reduced diastolic blood pressure by 2,20
mmHg and heart rate by 1,25 bpm, and increased flow-mediated dilation by
1,02\%, while the pooled reduction in systolic blood pressure (3,55
mmHg) was not statistically significant {[}37{]}. Mechanistically, these
outcomes align with polyphenol-driven modulation of endothelial nitric
oxide bioavailability, reduced oxidative quenching of NO, and
attenuation of low-grade inflammation.

In meat systems, grape seed ingredients are most often introduced as
grape seed extract during the mixing stage, so the phenolic fraction
becomes uniformly distributed within the protein-fat matrix before
forming and subsequent processing (Table 3). Across product types, the
most reproducible impact is on oxidative stability, typically tracked by
TBARS. Mechanistically, this aligns with the chemistry of grape seed
polyphenols which can quench lipid peroxyl radicals (LOO•) as
chain-breaking antioxidants and chelate transition metals, thereby
reducing initiation and propagation of lipid peroxidation driven by heme
iron and non-heme catalysts. The numeric signals in the table are
consistent with that model. Even when results are reported as relative
changes, Zhou et al. still observed a 30,3\% TBARS reduction, which is
the same direction expected from phenolic-driven suppression of lipid
oxidation {[}44{]}.

Because lipid oxidation and pigment oxidation are coupled, grape seed
additives also frequently affect colour stability and premature browning
behaviour. Xin et al. reported that redness on the raw surface at day 10
increased from 8,62 (control) to 13,20 at 0,75 g/kg, while TBARS stayed
within 0-0,51 mg MDA/kg-a pattern consistent with lowering oxidative
pressure on myoglobin and slowing conversion toward metmyoglobin
(brown), particularly at the oxygen-exposed surface {[}41{]}. In
parallel, texture and water-binding outcomes can shift because
polyphenols can interact with myofibrillar proteins, potentially
altering protein-protein aggregation and water immobilisation during
drying or heating. That helps explain why Zhou et al. saw reduced
expressible water by 26,4\% alongside improved sensory scores {[}44{]}.

{\bfseries Conclusion.} Analysis of the literature allows us to construct a
convincing scientific rationale for the selection of raw materials and
functional components. Lamb, as a meat ingredient for roulade, is
particularly vulnerable to lipid and pigment oxidation due to the
release of heme catalysts. In this context, sea buckthorn and grape
seeds form a complementary antioxidant platform: sea buckthorn combines
phenolic compounds, ascorbic acid, carotenoids and organic acids,
potentially affecting the redox balance and colour stability, while
grape seeds provide `polyphenolic support' (proanthocyanidins) capable
of breaking peroxide oxidation chains and binding transition metal ions,
inhibiting the decomposition of hydroperoxides and the accumulation of
secondary oxidation products. A comparison of studies shows that the
most controllable and reproducible method of introducing both additives
is in extract form, as it ensures accurate dosing, uniform distribution
of bioactive substances in the protein-fat system, and minimises the
risk of undesirable effects of dispersed particles on the cut and
texture of roll products.

Based on this, subsequent studies are expected to show a decrease in
oxidative degradation markers, more stable colour parameters and a
reduction in the formation of rancid aroma tones during storage, as well
as the preservation of technological suitability due to a reduction in
losses during heat treatment and textural characteristics. Subsequent
studies should focus on optimising extract doses and key technological
conditions with mandatory control of the trade-offs between antioxidant
effect and organoleptic indicators. This approach will allow us to
quantitatively substantiate the functional viability of the lamb and
plant antioxidants composition and develop a reproducible technological
scheme for meat rolls with high nutritional value.
\end{multicols}

\begin{center}
{\bfseries References}
\end{center}

\begin{refs}
1. Uzakov Y. et al. Collagen Hydrolysate-Cranberry Mixture as a
Functional Additive in Sausages //Processes. -2025. Vol.13(10):3233. DOI
\href{https://doi.org/10.3390/pr13103233}{10.3390/pr13103233}.

2. Kenenbay G. Et al. Marbling and Meat Quality of Kazakh Finewool
Purebred and Suffolk× Finewool Crossbred Sheep on an Intensive Fattening
Diet//Processes. -2025.-Vol.13(9):2874.
\href{https://doi.org/10.3390/pr13092874}{DOI 10.3390/pr13092874}.

3. Ding W. et al. Meat of sheep: insights into mutton evaluation,
nutritive value, influential factors, and interventions //Agriculture.
-2024. Vol.14. (7):1060. DOI
\href{https://doi.org/10.3390/agriculture14071060}{10.3390/agriculture14071060}.

4. Olas B. Sea buckthorn as a source of important bioactive compounds in
cardiovascular diseases //Food and Chemical Toxicology.
-2016.-Vol.97.-P.199-204.
\href{https://doi.org/10.1016/j.fct.2016.09.008}{DOI
10.1016/j.fct.2016.09.008}.

5. Gâtlan A. M., Gutt G. Sea buckthorn in plant based diets. An
analytical approach of sea buckthorn fruits composition: Nutritional
value, applications, and health benefits //environmental research and
public health. -2021.-Vol.18(17):8986. ~DOI
\href{https://doi.org/10.3390/ijerph18178986}{10.3390/ijerph18178986}.

6. Bhutani M. et al. Valorization of Grape By-Products: Insights into
Sustainable Industrial and Nutraceutical Applications//Future Foods.
-2025.-Р.1-20. DOI
\href{https://doi.org/10.1016/j.fufo.2025.100710}{10.1016/j.fufo.2025.100710}.

7. Verma P. et al. Structural chemistry to therapeutic functionality: A
comprehensive review on proanthocyanidins //Biocatalysis and
Agricultural
Biotechnology.-2023.-Vol.5:102963.\href{https://doi.10.1016/j.bcab.2023.102963}{DOI
10.1016/j.bcab.2023.102963}.

8. Junkuszew A. et al. Chemical composition and fatty acid content in
lamb and adult sheep meat //Archives Animal Breeding. -2020. Vol.63(2).
-P.261-268. \href{https://doi.org/10.5194/aab-63-261-2020}{DOI
10.5194/aab-63-261-2020}.

9. Suliman G. M. et al. A comparative study of sheep breeds: fattening
performance, carcass characteristics, meat chemical composition and
quality attributes //Frontiers in veterinary science.
-2021.-Vol.8:647192.
\href{https://doi.org/10.3389/fvets.2021.647192}{DOI
10.3389/fvets.2021.647192}.

10. Sabbioni A. et al. Allometric coefficients for physical-chemical
parameters of meat in local sheep breed //Small ruminant research.
-2019. Vol.174.-P.141-147. DOI
\href{https://doi.org/10.1016/j.smallrumres.2019.04.001}{10.1016/j.smallrumres.2019.04.001}.

11. Gama K. V. M. F. et al. Fatty acid, chemical, and tissue composition
of meat comparing Santa Inês breed sheep and Boer crossbreed goats
submitted to different supplementation strategies //Tropical animal
health and production. -2019.-Vol.52(2). -P.601-610.
\href{https://doi.org/10.1007/s11250-019-02047-1}{DOI
10.1007/s11250-019-02047-1}.

12. Erasmus S. W., Muller M., Hoffman L. C. Authentic sheep meat in the
European Union: Factors influencing and validating its unique meat
quality// Science of Food and Agriculture. -2017. -Vol.97(7). -
P.1979-1996. DOI 10.1002/jsfa.8180.

13. Geletu U. S. et al. Quality of cattle meat and its compositional
constituents //Veterinary medicine international. -2021.-Vol.2021(1).
DOI 10.1155/2021/7340495.

14. Picard B., Gagaoua M. Muscle fiber properties in cattle and their
relationships with meat qualities: An overview//Journal of Agricultural
and Food Chemistry. -2020.-Vol.68(22). -P.6021--6039. DOI
10.1021/acs.jafc.0c02086.

15. Venkata Reddy B. et al. Beef quality traits of heifer in comparison
with steer, bull and cow at various feeding environments//Animal Sci J.-
2015. -Vol.86(1). -P.1-16. DOI 10.1111/asj.12266.

16. Vicente F., Pereira P. C. Pork meat composition and health: A review
of the evidence//Foods. -2024.-Vol.13(12): 1905. DOI
10.3390/foods13121905.

17. Škrlep M. et al. The use of pork from entire male and
immunocastrated pigs for meat products-An overview with
recommendations//Animals. -2020.-Vol.10(10):1754. DOI
\href{https://doi.org/10.3390/ani10101754}{10.3390/ani10101754}

18. Aaslyng M. D., Meinert L. Meat flavour in pork and beef--From animal
to meal//Meat science. -2017.-Vol.132.-P.112-117.
\href{https://doi.org/10.1016/j.meatsci.2017.04.012}{DOI
10.1016/j.meatsci.2017.04.012}.

19. Liao J. et al. New insights into the effects of dietary amino acid
composition on meat quality in pigs: A review//Meat
science.-2025.-Vol.221:109721.
\href{https://doi.org/10.1016/j.meatsci.2024.109721}{DOI
10.1016/j.meatsci.2024.109721}.

20. Tkacz K. et al. Anti-oxidant and anti-enzymatic activities of sea
buckthorn (Hippophaë rhamnoides L.) fruits modulated by chemical
components//Antioxidants.-2019.-Vol.8(12):618.
\href{https://doi.org/10.3390/antiox8120618}{DOI 10.3390/antiox8120618}.

21. Bośko P. et al. Chemical composition and nutritive value of sea
buckthorn leaves//Molecules. -2024.Vol.29(15): 3550.
\href{https://doi.org/10.3390/molecules29153550}{DOI
10.3390/molecules29153550}.

22. Jaroszewska A., Biel W., Gibczyńska M. Analysis of the content of
macro-and microelements in the leaves of sea buckthorn (Hippophaë
rhamnoides L.). -2017.-Vol.338(44).-P.61-88.
\href{https://doi.org/10.21005/AAPZ2017.44.4.07}{DOI
10.21005/AAPZ2017.44.4.07}.

23. Criste A. et al. Phytochemical composition and biological activity
of berries and leaves from four Romanian sea buckthorn (Hippophae
rhamnoides L.) varieties //Molecules. -2020. Vol.25(5):1170.
\href{https://doi.org/10.3390/molecules25051170}{DOI
10.3390/molecules25051170}.

24. Zenkova M., Pinchykova J. Chemical composition of Sea-buckthorn and
Highbush Blueberry fruits grown in the Republic of Belarus //Food
Science and Applied Biotechnology. -2019. Vol.2(2). -P.121-129.
\href{https://doi.org/10.30721/fsab2019.v2.i2.59}{DOI
10.30721/fsab2019.v2.i2.59}.

25. Huang N. K. et al. Supplementation with Seabuckthorn Oil Augmented
in 16: 1n--7tIncreases SerumTrans-Palmitoleic Acid in Metabolically
Healthy Adults: A Randomized Crossover Dose-Escalation Study//The
Journal of Nutrition. -2020.-Vol.150(6). -P.1388-1396.
\href{https://doi.org/10.1093/jn/nxaa060}{DOI 10.1093/jn/nxaa060}.

26. Kim D. S., Iida F. Texture characteristics of sea buckthorn
(Hippophae rhamnoides) jelly for the elderly based on the gelling
agent//Foods.-2022.Vol.11(13):1892.
\href{https://doi.org/10.3390/foods11131892}{DOI 10.3390/foods11131892}.

27. Ren Z. et al. Effect of sea buckthorn on plasma glucose in
individuals with impaired glucose regulation: a two-stage randomized
crossover intervention study //Foods. -2021. Vol.10(4):804.
\href{https://doi.org/10.3390/foods10040804}{DOI 10.3390/foods10040804}.

28. Alimardanova M. K. et al. Effect of Hippophae rhamnoides Extract
Addition on the Quality and Safety of Traditional Kazakh Chunked
Delicacy ``Jaya''//Foods. -2025. Vol.14(21):3698.
\href{https://doi.org/10.3390/foods14213698}{DOI 10.3390/foods14213698}.

29. Bukarbayev K. et al. Effect of hemp protein and sea buckthorn
extract on quality and shelf life of cooked-smoked sausages//Foods.
-2025. Vol.14(15):2730. \href{https://doi.org/10.3390/foods14152730}{DOI
10.3390/foods14152730}.

30. Bobko M. Et al. Sensory properties of pork sausage after sea
buckthorn extract addition //Czech Journal of Food Sciences. -2025.
Vol.43(5).-P.320-325. \href{https://doi.org/10.17221/11/2025-CJFS}{DOI
10.17221/11/2025-CJFS}.

31. Wojtaszek A. et al. Physicochemical, Antioxidant, Organoleptic, and
Anti-Diabetic Properties of Innovative Beef Burgers Enriched with Juices
of Açaí (Euterpe oleracea Mart.) and Sea Buckthorn (Hippophae rhamnoides
L.) Berries //Foods. -2024.-Vol.13(19):3209.
\href{https://doi.org/10.3390/foods13193209}{DOI 10.3390/foods13193209}.

32. Guo Z. et al. Changes in chilled beef packaged in starch film
containing sea buckthorn pomace extract and quality changes in the film
during super-chilled storage //Meat Science. - 2021. Vol.182: 108620.
DOI
\href{https://doi.org/10.1016/j.meatsci.2021.108620}{10.1016/j.meatsci.2021.108620}.

33. Mesárošová A. et al. Effect of sea buckthorn (Hippophae rhamnoides
var. vitaminnaja) extract on spoilage of pork sausages//Journal of
microbiology, biotechnology and food sciences. -2024. -Vol.13(4):
e10420. DOI
\href{https://doi.org/10.55251/jmbfs.10420}{10.55251/jmbfs.10420}.

34. Baca-Bocanegra B. et al. Optimization of protein extraction of
oenological interest from grape seed meal using design of experiments
and response surface methodology//Foods. -2021. -Vol.10(1):79. DOI
\href{https://doi.org/10.3390/foods10010079}{10.3390/foods10010079}.

35. Carmona-Jiménez Y. et al. Fatty acid and tocopherol composition of
pomace and seed oil from five grape varieties southern Spain//Molecules.
-2022.-Vol.27(20):6980.DOI
\href{https://doi.org/10.3390/molecules27206980}{10.3390/molecules27206980}.

36. Ivanov Y., Atanasova M., Godjevargova T. Nutritional and Functional
Values of Grape Seed Flour and Extract for Production of Antioxidative
Dietary Supplements and Functional Foods//Molecules. - 2025.
Vol.30(9):2029.
\href{https://doi.org/10.3390/molecules30092029}{DOI10.3390/molecules30092029}.

37. Foshati S. et al. The effect of grape (Vitis vinifera) seed extract
supplementation on flow-mediated dilation, blood pressure, and heart
rate: A systematic review and meta-analysis of controlled trials with
duration-and dose-response analysis //Pharmacological research. -2022.
Vol.175:105905. \href{https://doi.org/10.1016/j.phrs.2021.105905}{DOI
10.1016/j.phrs.2021.105905}.

38. Pateiro M. et al. Oxidation stability of pig liver pâté with
increasing levels of natural antioxidants (grape and tea)//Antioxidants.
-2015. Vol.4(1). -P.102-123.
\href{https://doi.org/10.3390/antiox4010102}{DOI 10.3390/antiox4010102}.

39. Pateiro M. et al. Effect of addition of natural antioxidants on the
shelf-life of ``Chorizo'', a Spanish dry-cured sausage//Antioxidants.
-2015. Vol.4(1). -P.42-67.
\href{https://doi.org/10.3390/antiox4010042}{DOI 10.3390/antiox4010042}.

40. Gómez I. et al. Shelf life of ground beef enriched with omega‐3
and/or conjugated linoleic acid and use of grape seed extract to inhibit
lipid oxidation //Food Science \& Nutrition. -2016. Vol.4(1). -P.67-79.
\href{https://doi.org/10.1002/fsn3.251}{DOI 10.1002/fsn3.251}.

41. Xin L. U. O. et al. Effects of grape seed extract on meat color and
premature browning of meat patties in high-oxygen packaging //Journal of
Integrative Agriculture. -2022.-Vol.21(8).-P.2445-2455.
\href{https://doi.org/10.1016/S2095-3119(21)63854-6}{DOI
10.1016/S2095-3119(21)63854-6}.

42. Libera J. et al. Use of grape seed extract as a natural antioxidant
additive in dry-cured pork neck technology//Biotechnology and Food
Science. -2018. Vol.82(2).-P.141-150.
\href{https://doi.org/10.34658/bfs.2018.82.2.141-150}{DOI
10.34658/bfs.2018.82.2.141-150}.

43. Ivanov Y. et al. The Effect of Grape Seed Extract on Lipid
Oxidation, Color Change, and Microbial Growth in a Beef--Pork Sausage
Model System //Molecules. -2025.-Vol.30(8):1739.
\href{https://doi.org/10.3390/molecules30081739}{DOI
10.3390/molecules30081739}.

44. Zhou Y., Wang Q. Y., Wang S. Effects of rosemary extract, grape seed
extract and green tea polyphenol on the formation of N‐nitrosamines and
quality of western‐style smoked sausage //Journal of Food Processing and
Preservation.-2020.-Vol.44(6):e14459.DOI
\href{https://doi.org/10.1111/jfpp.14459}{10.1111/jfpp.14459}.
\end{refs}

\begin{info}
\hspace{1em}\emph{{\bfseries Information about the authors}}

Abdurasulova A.E. - PhD student, Master of Technical Sciences, Almaty
Technological University, Almaty, Kazakhstan, e-mail:
abdurasulova.aigerim@mail.ru;

Rskeldiyev B.A. - Doctor of Technical Sciences, Professor, Almaty
Technological University, Kazakhstan, Almaty; e-mail: berdan\_r@mail.ru.

\hspace{1em}\emph{{\bfseries Сведения об авторах}}

Абдурасулова А.Е. - магистр, докторант, Алматинский технологический
университет, Алматы, Казахстан, e-mail: abdurasulova.aigerim@mail.ru;

Рскелдиев Б.А. - д.т.н., профессор, Алматинский технологический
университет, Алматы, Казахстан, Алматы, e-mail: berdan\_r@mail.ru.
\end{info}
